\par\noindent\textbf{(1)} 令场包络 \(F(r)e^{i(l\phi+\beta z)}\)，径向方程 \(F''+(1/r)F'+[(n(r)k_0)^2-\beta^2-l^2/r^2]F=0\)。芯内 \(u^2=a^2(n_1^2k_0^2-\beta^2)>0\) 取 \(J_l(ur/a)\)，包层 \(w^2=a^2(\beta^2-n_2^2k_0^2)>0\) 取衰减的 \(K_l(wr/a)\)。
\par\noindent\textbf{(2)} 若近似标量场及径向导数在 \(r=a\) 连续，则 \(J_l'(u)u/J_l(u)=K_l'(w)w/K_l(w)\)。再联立 \(u^2+w^2=V^2\)，\(V=ak_0\sqrt{n_1^2-n_2^2}\)，求允许离散 \(\beta\) 值；精确矢量模式需考虑切向 \(E,H\) 共同连续。
\par\noindent\textbf{(3)} 当 \(wr/a\gg1\)，第二类修正Bessel函数 \(K_l(wr/a)\sim\sqrt{\pi a/(2wr)}e^{-wr/a}\)，因此包层强度主要按 \(e^{-2wr/a}\) 衰减，典型幅穿透长度 \(a/w=1/\sqrt{\beta^2-n_2^2k_0^2}\)。
\par\noindent\textbf{(4)} 弯曲使局部外侧相速度条件改变，包层原本倏逝的区域可能在远侧转为可辐射传播，导致弯曲损耗；模式接近截止时 \(w\to0\)、尾场变长，通常更易泄漏。单用直光纤的特征方程不能定量预测任意弯曲半径下的损耗。
